% ==========================================================
% CHƯƠNG 6: KẾT LUẬN VÀ HƯỚNG PHÁT TRIỂN
% ==========================================================

\chapter{Kết luận và hướng phát triển}
\label{ch:ket-luan-huong-phat-trien}

\section{Kết luận}
\label{sec:ket-luan}

Luận văn đã thực hiện thiết kế và kiểm chứng mô đun cấu hình thí nghiệm đa năng cho vệ tinh nhỏ theo hướng mô đun hóa. Hệ thống được tổ chức thành các khối chức năng rõ ràng, bao gồm khối xử lý trung tâm, khối giao tiếp, khối điều khiển ngoại vi, khối nguồn, khối điều khiển nhiệt độ và các thành phần phục vụ thí nghiệm quang học. Cách tiếp cận này giúp hệ thống có khả năng mở rộng, thuận lợi cho việc thay đổi cấu hình thí nghiệm và kiểm thử từng phần trước khi tích hợp toàn bộ.

Về phần cứng, luận văn đã phân tích yêu cầu hệ thống và triển khai thiết kế các board chức năng phục vụ mô đun cấu hình thí nghiệm. Các khối nguồn, giao tiếp, lưu trữ, cảm biến, điều khiển laser, điều khiển nhiệt và các đầu nối mở rộng được tổ chức nhằm đáp ứng nhu cầu vận hành của một cấu hình thí nghiệm. Thiết kế phần cứng không chỉ phục vụ một bài toán thử nghiệm cụ thể mà còn tạo nền tảng để tích hợp các ngoại vi khác trong những cấu hình sau.

Về phần mềm, luận văn đã xây dựng kiến trúc điều khiển gồm phần mềm trên máy tính, phần mềm trên MPU và firmware trên MCU. Phần mềm Mission Control hỗ trợ giám sát dữ liệu, cấu hình profile, điều khiển laser, quan sát đáp ứng nhiệt độ và quản lý trạng thái firmware. Firmware trên MCU đảm nhiệm các tác vụ thời gian thực như đọc cảm biến, điều khiển PID, xử lý profile nhiệt độ, điều khiển ngoại vi và phản hồi dữ liệu cho giao diện giám sát.

Một kết quả quan trọng khác là chức năng Bootloader đã được thiết kế và kiểm chứng. Bootloader cho phép cập nhật firmware cho MCU thông qua MPU bằng giao thức truyền dữ liệu có phản hồi ACK. Kết quả thử nghiệm cho thấy hệ thống có thể kết nối, truyền firmware theo từng gói, kiểm soát quá trình ghi dữ liệu và hỗ trợ chuyển sang ứng dụng chính sau khi cập nhật. Chức năng này có ý nghĩa thực tiễn đối với các hệ thống nhúng cần bảo trì phần mềm sau khi đã tích hợp vào môi trường vận hành.

Các kết quả trình bày ở Chương~\ref{ch:ket-qua} cho thấy hệ thống đã đạt được mục tiêu chính của luận văn: xây dựng được một nền tảng cấu hình thí nghiệm có khả năng điều khiển, giám sát, giao tiếp và cập nhật firmware. Giao diện điều khiển đã hỗ trợ quan sát trạng thái hệ thống theo thời gian thực; bộ điều khiển nhiệt độ có khả năng bám theo giá trị đặt; khối điều khiển laser đã được tích hợp vào giao diện vận hành; và Bootloader đã chứng minh được khả năng cập nhật chương trình thông qua MPU.

Nhìn chung, luận văn đã hoàn thành việc thiết kế nền tảng phần cứng và phần mềm cho mô đun thí nghiệm đa năng ở mức nguyên mẫu. Hệ thống đáp ứng được yêu cầu kiểm chứng chức năng, tạo cơ sở cho các bước thử nghiệm định lượng sâu hơn và phát triển thành một mô đun cấu hình thí nghiệm hoàn chỉnh hơn trong tương lai.

\section{Hạn chế}
\label{sec:han-che}

Bên cạnh các kết quả đạt được, hệ thống vẫn còn một số hạn chế cần được cải thiện. Thứ nhất, quá trình thử nghiệm hiện tại chủ yếu tập trung vào kiểm chứng chức năng và luồng vận hành, chưa thực hiện đầy đủ các đánh giá định lượng trong thời gian dài. Các thông số như độ ổn định nhiệt, độ lặp lại của đáp ứng điều khiển, sai số giữa các cảm biến và độ tin cậy khi vận hành liên tục cần được đo đạc với nhiều kịch bản hơn.

Thứ hai, chức năng điều khiển laser đã được tích hợp vào giao diện và có khả năng thao tác theo kênh, tuy nhiên cần bổ sung thêm các phép đo thực nghiệm để đánh giá độ chính xác công suất, độ đồng đều giữa các kênh và đáp ứng khi chuyển đổi nhanh giữa các trạng thái bật/tắt.

Thứ ba, Bootloader đã chứng minh được khả năng truyền và cập nhật firmware thông qua MPU, nhưng vẫn cần tiếp tục kiểm thử trong các điều kiện truyền thông không lý tưởng. Các tình huống như mất kết nối giữa chừng, sai dữ liệu, reset bất thường hoặc mất nguồn trong quá trình ghi Flash cần được mô phỏng và đánh giá để tăng độ tin cậy của cơ chế cập nhật.

Thứ tư, giao diện Mission Control hiện đã hỗ trợ các chức năng vận hành cơ bản, nhưng vẫn có thể được hoàn thiện hơn ở phần quản lý dữ liệu thí nghiệm, lưu log, xuất báo cáo kết quả và tự động hóa kịch bản kiểm thử. Những chức năng này sẽ giúp quá trình vận hành hệ thống thuận tiện và có tính lặp lại cao hơn.

\section{Hướng phát triển}
\label{sec:huong-phat-trien}

Trong thời gian tới, hệ thống có thể được phát triển theo các hướng sau:

\begin{itemize}
    \item \textbf{Đánh giá định lượng hệ thống điều khiển nhiệt:} Thực hiện các bài thử nghiệm dài hạn với nhiều mức nhiệt độ đặt khác nhau để đánh giá độ ổn định, thời gian đáp ứng, độ quá điều chỉnh và sai số xác lập của bộ điều khiển PID.

    \item \textbf{Hoàn thiện kiểm thử cụm laser:} Bổ sung phép đo công suất quang, độ đồng đều giữa các kênh laser và độ ổn định khi vận hành theo profile tự động. Kết quả này giúp đánh giá chính xác hơn khả năng phục vụ các thí nghiệm quang học.

    \item \textbf{Tăng độ tin cậy của Bootloader:} Bổ sung các cơ chế bảo vệ như xác thực firmware, kiểm tra chữ ký số, cơ chế rollback khi cập nhật thất bại và kiểm thử khả năng khôi phục sau lỗi nguồn hoặc lỗi truyền thông.

    \item \textbf{Mở rộng giao diện Mission Control:} Phát triển thêm chức năng lưu log tự động, xuất dữ liệu thí nghiệm, hiển thị cảnh báo lỗi, quản lý nhiều profile và hỗ trợ chạy kịch bản kiểm thử tự động.

    \item \textbf{Tích hợp thêm cảm biến và ngoại vi thí nghiệm:} Tận dụng kiến trúc mô đun để kết nối thêm các cảm biến hoặc cơ cấu chấp hành khác, từ đó mở rộng phạm vi ứng dụng của payload cho nhiều loại thí nghiệm.

    \item \textbf{Kiểm thử môi trường và độ bền phần cứng:} Thực hiện các bài kiểm tra rung, nhiệt, chu kỳ nguồn và vận hành liên tục nhằm đánh giá khả năng đáp ứng của hệ thống trong các điều kiện gần với môi trường làm việc thực tế.

    \item \textbf{Chuẩn hóa quy trình tích hợp hệ thống:} Xây dựng quy trình kiểm thử từng board, kiểm thử liên board và kiểm thử toàn hệ thống để giảm rủi ro trong các lần tích hợp sau.
\end{itemize}

Các hướng phát triển trên sẽ giúp mô đun cấu hình thí nghiệm tiến gần hơn đến một nền tảng hoàn chỉnh, có thể tái sử dụng cho nhiều nhiệm vụ khác nhau trên vệ tinh nhỏ. Đặc biệt, việc kết hợp giữa kiến trúc phần cứng mô đun, firmware có khả năng cập nhật và giao diện vận hành tập trung sẽ tạo điều kiện thuận lợi cho quá trình mở rộng hệ thống trong các nghiên cứu tiếp theo.
